% ===================================================================
%  CADRO N.º 5 — A casa e a súa construción.
%  Traduction 2026 d'apres texto/io, controlee sur texto/fr, texto/es
%  et texto/pt. Consigne : voir texto/gl/10-cadro-01.tex.
%
%  LE TABLEAU DU BATIMENT EST UN DIALOGUE, et c'est le seul du
%  volume. Le galicien y demande une decision que ni le portugais ni
%  l'espagnol n'imposent de la meme facon : la place du pronom.
%  L'affirmative postpose — « taspeñarei », « cultívaas », « dálle » —
%  mais la negative, l'interrogative, la subordonnee et un mot comme
%  « xa » ou « ben » anteposent : « non o vexo », « que lle dá a
%  man ». Le portugais fait de meme, l'espagnol antepose partout. Une
%  colonne galicienne ecrite sur l'espagnol dirait « lo veo » ou il
%  faut « véxoo », et rien de mecanique ne l'attraperait.
%
%  LES METIERS DU BATIMENT SONT PLUS PROCHES DE L'ESPAGNOL QUE CEUX
%  DU CORPS, et c'est encore le lieu de l'apprentissage qui le dit :
%  le magon, le charpentier, le serrurier s'apprennent sur un
%  chantier, et le chantier est le meme des deux cotes de la
%  frontiere. Restent quelques mots que le galicien seul a :
%  « albanel » le magon, contre « albañil » et « pedreiro » ;
%  « lousa » l'ardoise, qui est aussi l'ardoise de l'ecolier du
%  tableau 1 ; « fiestra » la fenetre, a cote de « ventá ».
% ===================================================================

%%K t05-apar-1 apar
\VUsaut{6.0mm}
\VUtitre{15.0pt}{60}{CADRO N.º 5}
\VUsaut{1.4mm}
\VUtitre{11.4pt}{0}{A casa e a súa construción.}
\VUsaut{2.2mm}

%%K t05-tit-1 sub
\VUpk{10.2pt}{40}{Bo encontro.}
\VUsaut{3.0mm}

%%K t05-01-1 p
\textsc{Xoán}. --- Meu tío, gustaríame moito saber como se constrúe unha
\VUgras{casa}.

%%K t05-01-2 p
\textsc{O tío} \textsuperscript{(80)}. --- Iso é moi doado, amigo meu; vén comigo e
amosareiche a que o señor Bernardo \textsuperscript{(79)} manda construír e na que eu vivirei.

%%K t05-01-3 p
\textsc{Xoán}. --- E quen debuxou o \VUgras{plano} \textsuperscript{(76)} desta casa?

%%K t05-01-4 p
\textsc{O tío}. --- O \VUgras{arquitecto} \textsuperscript{(75)}; presentou un debuxo moi
claro que foi aprobado, e o \VUgras{contratista} \textsuperscript{(77)} comezou os traballos.

%%K t05-01-5 p
\textsc{Xoán}. --- E quen é o contratista?

%%K t05-01-6 p
\textsc{O tío}. --- É o señor Blanc, o pai do teu amigo Xacobe. Dirixe os obreiros xunto co
señor Roux, o arquitecto.

%%K t05-01-7 p
Vaia! Xusto agora chega o señor Blanc coa súa \VUgras{regra} \textsuperscript{(78)}, e o señor
Roux co seu plano.

%%K t05-01-8 p
Bos días, señores. Como están?

%%K t05-01-9 p
\textsc{O señor Roux}. --- Moi ben, grazas. Bos días, Xoán; canto medrou este neno!

%%K t05-01-10 p
\textsc{O tío}. --- Si, de certo, é un homiño. É moi serio; hai que o informar de todo canto
ve. Hoxe quixera saber como se constrúe unha casa.

%%K t05-tit-2 sub
\VUpk{10.2pt}{40}{Os obreiros.}
\VUsaut{3.0mm}

%%K t05-01-11 p
\textsc{O señor Blanc}. --- Pois ven comigo, amiguiño, que che explico axiña, que quero moito
os rapaces que se queren instruír.

%%K t05-01-12 p
Ves ese obreiro que ten unha \VUgras{pa} \textsuperscript{(82)} na man: é un
\VUgras{peón de terra} \textsuperscript{(81)}. Escava unha \VUgras{gabia}
\textsuperscript{(46)} para poñer a canalización da auga. Tamén escavou o chan no sitio onde
está a casa, para que o albanel erguese os catro \VUgras{muros}. A parte deses muros que está
no chan chámase os \VUgras{cimentos} \textsuperscript{(2)}. O espazo que fica entre os
cimentos forma o \VUgras{soto} \textsuperscript{(3)}. Ese soto ten unha fiestra pequena
chamada \VUgras{tragaluz} \textsuperscript{(5)}, unha
\VUgras{porta} \textsuperscript{(4)} e unha escaleira.

%%K t05-01-13 p
O \VUgras{albanel} \textsuperscript{(108)} seguiu erguendo os muros deixando aberturas para as
\VUgras{portas} \textsuperscript{(8)} e as \VUgras{fiestras} \textsuperscript{(17)}.

%%K t05-01-14 p
\textsc{Xoán}. --- Pero ese albanel non o vexo!

%%K t05-01-15 p
\textsc{O señor Blanc}. --- Agora rematou de traballar na casa. Está preto da
\VUgras{cabaleriza} \textsuperscript{(54)}, no seu \VUgras{andamio}
\textsuperscript{(110)}; constrúe o muro do \VUgras{lavadoiro} \textsuperscript{(56)}.

%%K t05-01-16 p
Ten unha trulla, unha artesa e unha \VUgras{prumada} \textsuperscript{(109)}. Pero ti non
lembrarás eses nomes.

%%K t05-01-17 p
\textsc{Xoán}. --- Claro que si, que axiña lle pedirei ao meu tío unha trulliña e unha artesa,
e eu mesmo farei unha prumada cun cordel.

%%K t05-tit-3 sub
\VUsaut{3.0mm}
\VUpk{10.2pt}{40}{Os materiais.}
\VUsaut{3.0mm}

%%K t05-01-18 p
\textsc{O tío}. --- Ha! Moi ben. Pero dime, Xoán, que fai falta para construír unha casa?

%%K t05-01-19 p
\textsc{Xoán}. --- Fan falta \VUgras{pedras de cantaría} \textsuperscript{(59)} e
\VUgras{morteiro} \textsuperscript{(65)}.

%%K t05-01-20 p
\textsc{O tío}. --- Xusto. Mira ese home do sombreiro grande, no seu \VUgras{carro}
\textsuperscript{(136)}. É o \VUgras{carreteiro} \textsuperscript{(135)}. Cada día trae aquí
\VUgras{bloques de pedra} \textsuperscript{(60)}. Descarga o seu carro no patio, e os
\VUgras{canteiros} \textsuperscript{(83)} labran os bloques e sérranos coa súa
\VUgras{serra de pedra} \textsuperscript{(85)}. Un \VUgras{peón} \textsuperscript{(146)}
lévaos ao albanel, que os coloca.

%%K t05-01-21 p
Mentres tanto, o \VUgras{amasador} \textsuperscript{(87)} prepara o morteiro. Precisa de
\VUgras{area} \textsuperscript{(62)}, \VUgras{cal} \textsuperscript{(63)} e
\VUgras{auga} \textsuperscript{(64)}. A cal está preto del nun \VUgras{saco}
\textsuperscript{(71)}; un \VUgras{servente de albanel} \textsuperscript{(111)} tráelle a area
nunha \VUgras{carretilla} \textsuperscript{(112)}, e o
\VUgras{aprendiz} \textsuperscript{(113)} está na bomba (58). Enche un \VUgras{balde}
\textsuperscript{(114)}. O morteiro está axiña listo. O \VUgras{portamorteiro}
\textsuperscript{(107)} cólleo e lévao arriba, pola escaleira, ao andamio, onde traballa un
albanel que remata ese lado da casa.

%%K t05-01-22 p
E agora, como se chama o obreiro que traballa o \VUgras{tellado} \textsuperscript{(31)}?

%%K t05-01-23 p
\textsc{Xoán}. --- É o \VUgras{tellador} \textsuperscript{(118)}.

%%K t05-tit-4 sub
\VUsaut{3.0mm}
\VUpk{10.2pt}{40}{O tellado da casa.}
\VUsaut{3.0mm}

%%K t05-01-24 p
\textsc{O señor Blanc}. --- Si, pero primeiro vén o \VUgras{carpinteiro}
\textsuperscript{(115)}.

%%K t05-01-25 p
O carpinteiro traballa o \VUgras{armazón} \textsuperscript{(35)}. Xunta as
\VUgras{trabes} \textsuperscript{(37)} con \VUgras{caravillas}
\textsuperscript{(117)}. Eses obreiros de aí arriba, coas súas mangas anchas de veludo, son
carpinteiros. Un ten unha traberiña, e o outro corta unha caravilla coa súa
\VUgras{machada} \textsuperscript{(116)}. O que está preto da \VUgras{cheminea}
\textsuperscript{(41)} é o tellador. Crava as ripas nas (*) vigas, e as
\VUgras{lousas} \textsuperscript{(66)} nas ripas. Ás veces pon tellas.

%%K t05-noto-1 noto
\VUnotes{143.2mm}{%
(*) \VUgras{Balk-o} = Al. \textit{Balken}, Ing. \textit{joist}, Fr. \textit{solive}, It.
\textit{travicello}, Ru. \textit{balka}, Gal., Esp. e Port. \textit{viga}. Raíz técnica aínda
non adoptada, pero que cómpre propoñer para traducir o sentido do francés \textit{solive}, que
non é nin \textit{trabo} (Fr. \textit{poutre}) nin unha viguiña. É un pau de madeira
escuadrado que sostén o piso e o teito.%
}

%%K t05-01-26 p
O tellado da cabaleriza está \VUgras{cuberto de tellas} \textsuperscript{(55)}, o das casas
está \VUgras{cuberto de lousa} \textsuperscript{(32)} e o do alpendre é de
\VUgras{cinc} \textsuperscript{(24)}.

%%K t05-01-27 p
\textsc{Xoán}. --- E o tellador traballa tamén os tellados de cinc?

%%K t05-01-28 p
\textsc{O señor Blanc}. --- Non, iso faino o \VUgras{cincador} \textsuperscript{(121)} ou o
\VUgras{chumbeiro}, que é o que pon a \VUgras{claraboia} \textsuperscript{(38)} no tellado da
casa. Pon tamén as \VUgras{canles da chuvia} \textsuperscript{(43)} e as
\VUgras{canles do tellado} \textsuperscript{(42)}. Solda xunto o cinc e a folla de lata. Fixou
o \VUgras{pararraios} \textsuperscript{(40)} e a \VUgras{veleta} \textsuperscript{(39)} no
\VUgras{fastial} \textsuperscript{(36)}.

%%K t05-01-29 p
\textsc{Xoán}. --- E así a casa está rematada?

%%K t05-tit-5 sub
\VUsaut{3.0mm}
\VUpk{10.2pt}{40}{As ferramentas.}
\VUsaut{3.0mm}

%%K t05-01-30 p
\textsc{O señor Blanc}. --- Non de todo. O \VUgras{xeseiro} \textsuperscript{(99)} constrúe
con \VUgras{ladrillos} \textsuperscript{(61)} os tabiques que a dividen en distintos cuartos.
Recobre con \VUgras{xeso} \textsuperscript{(70)} os
\VUgras{teitos} \textsuperscript{(26)} e mais os muros. Agora traballa o teito do
\VUgras{alpendre} \textsuperscript{(33)}. Ten, coma o albanel, unha \VUgras{trulla}
\textsuperscript{(100)} e unha \VUgras{artesa} \textsuperscript{(101)}.

%%K t05-01-31 p
Despois o ebanista traballa as portas, as fiestras, os \VUgras{parqués} \textsuperscript{(16)}
e as \VUgras{contras} \textsuperscript{(19)}.

%%K t05-01-32 p
Agora traballa no patio, no seu \VUgras{banco de traballo} \textsuperscript{(90)}. Ten unha
\VUgras{serra} \textsuperscript{(94)} para serrar a madeira (69), unha
\VUgras{garlopa} \textsuperscript{(91)}, un \VUgras{gato} \textsuperscript{(92)}, un
\VUgras{formón} \textsuperscript{(94 bis)}, un \VUgras{compás}
\textsuperscript{(95)} e un \VUgras{mazo} de madeira \textsuperscript{(95 bis)}.

%%K t05-01-33 p
\textsc{Xoán}. --- Ha! Pero é máis divertido facer de ebanista ca de albanel. Meu tío, ti
hasme mercar un banco de traballo, unha serra e unha garlopa, que axiña farei unha caseta para
Medor.

%%K t05-01-34 p
\textsc{O tío}. --- Si, rapaciño.

%%K t05-01-35 p
\textsc{Xoán}. --- Que contento estou! E quen é ese que está de pé preto da porta de entrada?

%%K t05-tit-6 sub
\VUsaut{3.0mm}
\VUpk{10.2pt}{40}{A pintura.}
\VUsaut{3.0mm}

%%K t05-01-36 p
\textsc{O señor Blanc}. --- É o \VUgras{pintor} \textsuperscript{(102)}. Está de pé na súa
escaleira cun pote de \VUgras{pintura} \textsuperscript{(103)} e o seu
\VUgras{pincel} \textsuperscript{(104)}. Pinta a madeira da porta de entrada para que
se conserve máis tempo.

%%K t05-01-37 p
Tamén é el quen pega os papeis nas vivendas.

%%K t05-01-38 p
O \VUgras{tapiceiro} \textsuperscript{(107)}, que está nunha \VUgras{escaleira de man}
\textsuperscript{(106)}, no \VUgras{balcón} \textsuperscript{(28)}, pon un
\VUgras{estor} \textsuperscript{(29)}.

%%K t05-01-39 p
O obreiro que está de pé preto da fiestra grande é o \VUgras{vidreiro} \textsuperscript{(96)}.
Fixa os \VUgras{vidros} \textsuperscript{(18)} con
\VUgras{macilla} \textsuperscript{(73)}. Ese outro obreiro, axeonllado preto da porta
do soto, é o \VUgras{serralleiro} \textsuperscript{(97)}. Pon unha
\VUgras{pechadura} \textsuperscript{(6)}. A súa \VUgras{lima}
\textsuperscript{(98)} está preto del.

%%K t05-01-40 p
\textsc{Xoán}. --- E tamén é serralleiro ese obreiro que bate co seu
\VUgras{martelo} \textsuperscript{(84)} nun ferro?

%%K t05-01-41 p
\textsc{O señor Blanc}. --- Para dicilo mellor, é ferreiro (125). Forxa
\VUgras{ferraxes} \textsuperscript{(67)} para o \VUgras{electricista}
\textsuperscript{(124)}, que está no tellado da \VUgras{cocheira} \textsuperscript{(51)}. Ten
unha \VUgras{forxa} \textsuperscript{(126)}, unha
\VUgras{zafra} \textsuperscript{(127)} e unhas \VUgras{tenaces}
\textsuperscript{(128)}.

%%K t05-01-42 p
O electricista fixa o \VUgras{fío} de \VUgras{cobre} \textsuperscript{(44)} do teléfono.

%%K t05-tit-7 sub
\VUpk{10.2pt}{40}{A horta.}
\VUsaut{3.0mm}

%%K t05-01-43 p
\textsc{O tío}. --- Ao fondo do \VUgras{patio} \textsuperscript{(45)}, preto da
\VUgras{reixa} \textsuperscript{(47)} e do \VUgras{portalón}
\textsuperscript{(48)}, ves o \VUgras{xardineiro} \textsuperscript{(129)}. Prepara a
\VUgras{hortiña} \textsuperscript{(49)}. Cavou a terra coa súa \VUgras{pa}
\textsuperscript{(131)}, sementa sementes e rastra co \VUgras{rastro} \textsuperscript{(130)}.
Despois, coa súa \VUgras{regadeira} \textsuperscript{(132)}, regará os \VUgras{canteiros}
\textsuperscript{(150)} e as plantas irán medrando.

%%K t05-01-44 p
O \VUgras{cabaleirizo} \textsuperscript{(133)}, que acaba de coidar o cabalo e de lle dar a
comida, limpa o \VUgras{coche} \textsuperscript{(52)} cunha esponxa, unha
\VUgras{bomba de man} \textsuperscript{(134)} e un balde de auga. Despois volverá
metelo na cocheira e pechará a porta co groso \VUgras{ferrollo} \textsuperscript{(53)}.

%%K t05-01-45 p
Coido que xa estás contento; tiveches explicacións dabondo.

%%K t05-01-46 p
\textsc{Xoán}. --- Si, tío meu, pero de boa gana vería o interior da casa.

%%K t05-01-47 p
\textsc{O tío}. --- Iso será mañá. O señor Roux explicarache o plano e dirache o nome de cada
cuarto.

%%K t05-tit-8 sub
\VUsaut{3.0mm}
\VUpk{10.2pt}{40}{Distribución interior da casa.}
\VUsaut{2.4mm}

%%K t05-apar-2 apar
{\centering\textit{(Véxase o plano.)}\par}
\VUsaut{2.4mm}

%%K t05-01-48 p
\textsc{O arquitecto}. --- Ven, Xoán, que che vou amosar as distintas partes da casa.

%%K t05-01-49 p
Entremos no \VUgras{andar baixo} \textsuperscript{(7)}. Este é o botón do
\VUgras{timbre} \textsuperscript{(11)}. Subimos os tres \VUgras{chanzos}
\textsuperscript{(14)} do \VUgras{perrón} \textsuperscript{(9)}, cruzamos o
\VUgras{limiar} \textsuperscript{(10)} da \VUgras{porta de entrada}
\textsuperscript{(8)} e xa estamos ao pé da \VUgras{escaleira} \textsuperscript{(13)}.

%%K t05-01-50 p
\textsc{Xoán}. --- Iso é o corredor.

%%K t05-01-51 p
\textsc{O arquitecto}. --- Non, iso é o \VUgras{vestíbulo} \textsuperscript{(12)}. O
\VUgras{corredor} \textsuperscript{(a)} está na punta. Á dereita están o
\VUgras{salón} \textsuperscript{(b)} e o \VUgras{comedor} \textsuperscript{(c)};
fronte ao comedor están a \VUgras{cociña} \textsuperscript{(d)} e o \VUgras{cuarto do servizo}
\textsuperscript{(e)}, e despois, diante, o \VUgras{despacho} \textsuperscript{(f)}. O
\VUgras{alpendre} \textsuperscript{(23)} coa súa
\VUgras{porta de vidro} \textsuperscript{(25)} está preto do salón.

%%K t05-01-52 p
Imos subir agora a escaleira. Agárrate ao \VUgras{pasamáns} \textsuperscript{(15)} e conta os
chanzos. Son catorce. Xa estamos no \VUgras{rechán} \textsuperscript{(m)}: é o primeiro
\VUgras{andar} \textsuperscript{(27)}.

%%K t05-tit-9 sub
\VUsaut{3.0mm}
\VUpk{10.2pt}{40}{O primeiro andar.}
\VUsaut{3.0mm}

%%K t05-01-53 p
Fronte á escaleira están o \VUgras{retrete} \textsuperscript{(20)} e o \VUgras{cuarto de aseo}
\textsuperscript{(j)}. Á dereita un belo \VUgras{dormitorio} \textsuperscript{(g)} con dúas
fiestras na \VUgras{fachada} \textsuperscript{(1)} da casa. Á esquerda están o
\VUgras{cuarto de baño} \textsuperscript{(1)} e o
\VUgras{cuarto de hóspedes} \textsuperscript{(i)} cunha grande \VUgras{alcoba}
\textsuperscript{(h)}: preto do dormitorio está o \VUgras{cuarto dos nenos}
\textsuperscript{(k)}. Dá a unha ampla \VUgras{terraza} \textsuperscript{(21)}. Debaixo desa
terraza hai outro retrete (20) e, na parede, pende a \VUgras{campá} \textsuperscript{(22)} que
toca para a cea.

%%K t05-01-54 p
\textsc{Xoán}. --- Imos ao \VUgras{segundo andar}.

%%K t05-01-55 p
\textsc{O arquitecto}. --- Non hai segundo andar. Baixo o tellado só hai
\VUgras{faiados} \textsuperscript{(33)} e o faiado grande (34). Rematamos a visita;
podemos baixar.

%%K t05-01-56 p
\textsc{Xoán}. --- Grazas, señor; foi moi interesante. Axiña lle vou contar a meu pai todo
canto vin.
